\chapter{Lecture 2: 이동 방식 (Locomotion Concepts)}
\label{ch:lec02}

% cleveref names follow the language of the surrounding prose (hw2.tex 4-14).
\crefname{equation}{식}{식}
\Crefname{equation}{식}{식}
\crefname{figure}{그림}{그림}
\Crefname{figure}{그림}{그림}
\crefname{table}{표}{표}
\Crefname{table}{표}{표}
\crefname{section}{절}{절}
\Crefname{section}{절}{절}
\crefname{chapter}{장}{장}
\Crefname{chapter}{장}{장}

\paragraph{In brief.}
\begin{description}
  \item[Why] 로봇이 세상을 움직여 다니는 방법(이동, locomotion)을 결정해야 실제로
        환경 안에서 일을 할 수 있다. 다리로 걸을지 바퀴로 구를지, 몇 개의 관절과
        몇 가지 걸음새(gait)를 쓸지는 에너지·제어 복잡도·지형 적응성을 모두 바꾼다.
  \item[When] 로봇이 지면(또는 물, 공기)과 물리적으로 접촉해 움직이는 모든 상황.
        바퀴형 이동 로봇, 다리 로봇, 비행 로봇, 수중 로봇이 전부 이 범주에 든다.
  \item[Where] see-think-act 순환의 마지막 단계인 행동(acting) 블록, 그중에서도
        구동기 명령(actuator commands)이 실제 움직임으로 바뀌는 자리에 해당한다.
  \item[How] 이동 방식을 안정성(stability), 접촉 특성(contact characteristics),
        환경 종류(type of environment)로 특성화하고, 위치와 방향을 벡터와 회전
        행렬(rotation matrix)로 표현해 정량적으로 다룬다.
\end{description}

\section{배경: see-think-act에서 acting이 하는 일}
\label{sec:lec02-bg}

이 절의 출처는 \emph{Lecture 2: Locomotion Concepts} 슬라이드 1--6쪽이다.

\subsection{이동(locomotion)이란 무엇인가}
\label{sec:lec02-bg-def}

이동(locomotion)은 로봇(vehicle)과 그 환경 사이의 물리적 상호작용이다. 조금 더
풀면, 로봇이 앞으로 나아가려면 지면이나 물, 공기와 힘을 주고받아야 하고, 그 힘을
만들어 내는 것이 구동기(actuator)와 그것을 둘러싼 기구(mechanism)다.
% src: Lecture2-Locomotion.pdf slide 6
그래서 이동을 이해한다는 것은 "로봇이 어디로 가는가"가 아니라 "로봇이 그 힘을
어떻게 만들어 내는가"를 이해하는 일이다.

이 장은 \Cref{ch:lec01}에서 그린 see-think-act 순환(\Cref{fig:lec01-see-think-act})의
마지막 조각을 채운다. 그 순환은 감지(sensing) $\to$ 인지(cognition) $\to$
행동(acting)의 흐름이었고, 이동은 그 마지막 화살표, 즉 구동기 명령(actuator
commands)이 실제 세계의 움직임으로 바뀌는 실행(execution) 블록에 해당한다.
% src: Lecture2-Locomotion.pdf slide 2
앞의 두 블록(감지, 인지)이 "어디에 있고 어디로 가야 하는가"를 계산했다면, 이
장은 "그 결정을 물리적으로 어떻게 실현하는가"를 다룬다.

\subsection{선행 개념 네 가지}
\label{sec:lec02-bg-prereq}

본문에서 바로 쓰이는 용어를 먼저 정리한다.

\begin{description}
  \item[자유도(degree of freedom, DOF)] 하나의 강체(rigid body)나 관절이 독립적으로
        움직일 수 있는 방향의 개수. 다리 하나에 DOF가 많을수록 더 다양한 자세를
        만들 수 있지만, 그만큼 제어해야 할 변수도 늘어난다.
  \item[무게중심(center of gravity, COG)] 로봇의 질량이 평균적으로 쏠려 있는
        지점. 다리 로봇이 넘어지는지 아닌지는 결국 COG가 지지 다각형(support
        polygon, 접촉점들이 만드는 다각형) 안에 있는지로 판가름 난다.
  \item[정적 안정성과 동적 안정성(static vs.\ dynamic stability)] 정적으로
        안정하다는 것은 로봇이 그 순간 멈춰도 넘어지지 않는다는 뜻이고, 동적으로만
        안정하다는 것은 계속 움직여야만 넘어지지 않는다는 뜻이다. 자전거가 서
        있을 때는 넘어지지만 달릴 때는 넘어지지 않는 것과 같은 이치다. 자세한
        대비는 \Cref{sec:lec02-static-dynamic}에서 다룬다.
  \item[걸음새(gait)] 다리 여러 개를 가진 로봇이 다리를 드는 시점(lift)과
        내려놓는 시점(touch-down)을 어떤 순서로 반복하는지를 정한 패턴. 사람이
        걷는 방식과 뛰는 방식이 서로 다른 gait인 것과 같다. \Cref{sec:lec02-gaits}에서
        정량적으로 다룬다.
\end{description}

\section{왜 필요한가: 걸을 것인가, 구를 것인가}
\label{sec:lec02-why}

\paragraph{무슨 문제가 있었나.}
자연은 이동 방법을 아주 많이 진화시켰지만, 그중 바퀴를 쓰는 경우는 거의 없다.
반면 사람이 만든 기술 시스템은 대부분 바퀴나 무한궤도(caterpillar)를 쓴다.
% src: Lecture2-Locomotion.pdf slide 4
왜 이런 차이가 생기는가. 바퀴가 회전하려면 몸체와 분리된 채 계속 미끄러지듯 도는
관절이 필요한데, 생물학적 진화는 그런 회전 관절(rotating joint)을 만들어 내지
못했다. 반대로 다리로 걷는 방식은 자연에서는 흔하지만, 공학적으로는 구동기 수가
많고 제어가 훨씬 어렵다.

\paragraph{그래서 무엇을 바꿨나.}
그래서 로봇을 설계할 때는 "걸을 것인가, 구를 것인가(walking or rolling?)"라는
질문에 명시적으로 답해야 한다. 슬라이드가 제시하는 기준은 네 가지다:
구동기(actuator) 개수, 구조적 복잡도(structural complexity), 제어 비용(control
expense), 에너지 효율(energy efficiency), 그리고 굴러야 할 지형(terrain, 평지·무른
땅·오르막)이다.
% src: Lecture2-Locomotion.pdf slide 5

\paragraph{왜 그 변경이 문제를 푸는가.}
이 네 기준 중 에너지 효율은 숫자로 비교할 수 있다는 점에서 특히 유용하다. 그
숫자가 수송 비용(cost of transportation, CoT)이고, \Cref{sec:lec02-cot}에서
바로 다룬다. 나머지 세 기준(구동기 수, 구조 복잡도, 제어 비용)은 다리 로봇 쪽이
불리한 경우가 많다는 점을 \Cref{sec:lec02-walking-robots}에서 다시 확인한다.

\section{무엇인가}
\label{sec:lec02-what}

\subsection{걸을 것인가 구를 것인가: 수송 비용}
\label{sec:lec02-cot}

\paragraph{In brief.}
\begin{description}
  \item[Why] 걷기와 구르기 중 어느 쪽이 에너지를 덜 쓰는지 지형·속도별로 비교하려면
        질량과 이동 거리를 정규화한 공통 척도가 필요하다.
  \item[When] 이동 방식을 에너지 관점에서 비교할 때, 또는 로봇의 이동 효율을
        생물(동물)과 비교할 때 쓴다.
  \item[Where] 로봇 몸체 전체의 에너지 소비를 다룬다. 개별 관절의 토크나 전류
        같은 세부 사항은 다루지 않는다.
  \item[How] 소비한 에너지(또는 일률)를 무게와 이동 거리(또는 속도)로 나눠
        무차원 값을 만든다.
\end{description}

로봇 두 대를 비교한다고 하자. 하나는 무겁고 하나는 가볍다면, 단순히 "에너지를
얼마나 썼는가"만으로는 비교가 불공평하다. 무거운 로봇은 원래 더 많은 에너지가
필요하기 때문이다. 그래서 에너지를 무게로 나누어 표준화한 값이 필요하고, 그것이
수송 비용(cost of transportation, CoT)이다.
% src: Lecture2-Locomotion.pdf slide 5

\begin{equation}
\label{eq:lec02-cot}
\mathrm{CoT} = \frac{E}{m\,g\,d} = \frac{P}{m\,g\,v}.
\end{equation}

기호는 등장 순서대로 다음을 뜻한다.

\begin{description}
  \item[$E$] 이동에 사용한 에너지(energy), 단위는 줄(J).
  \item[$m$] 로봇의 질량(mass), 단위는 kg.
  \item[$g$] 중력가속도(gravity), 지구 표면에서 약 $9.8\,\text{m/s}^2$.
  \item[$d$] 이동한 거리(distance), 단위는 m.
  \item[$P$] 순간 일률(power), 단위는 와트(W). $E$를 이동에 걸린 시간으로 나눈 값이다.
  \item[$v$] 속도(speed), 단위는 m/s. $d$를 시간으로 나눈 값이므로 오른쪽 등식은
        왼쪽 등식의 분자·분모를 모두 시간으로 나눈 것과 같다.
\end{description}
% src: Lecture2-Locomotion.pdf slide 5

$m g$는 로봇의 무게(weight, 단위 N)이므로, CoT는 "단위 무게를 단위 거리만큼
옮기는 데 쓴 에너지"다. 분자(J)와 분모($\text{N}\cdot\text{m} = \text{J}$)의
단위가 같으므로 CoT 자체는 \emph{무차원(dimensionless)} 값이다. CoT가 작을수록
같은 무게를 같은 거리만큼 옮기는 데 에너지를 덜 쓴다는 뜻이니 효율이 좋다.

걷기와 뛰기는 여기에 추가 손실이 있다. 다리로 이동할 때는 몸통, 즉 무게중심
(COG, \Cref{sec:lec02-bg-prereq})이 한 걸음마다 위아래로 오르내린다. 그 오르내림
자체가 일이므로(중력을 거슬러 COG를 들어 올렸다가 다시 내려놓는 일), 순수하게
수평으로 미끄러지는 바퀴보다 다리 쪽이 같은 거리를 갈 때 여분의 에너지를 더
쓴다.
% src: Lecture2-Locomotion.pdf slide 5

\paragraph{예제 -- 수치 예시(저자가 구성함, 슬라이드 값 아님).}
질량 $m = 50\,\text{kg}$인 다리 로봇이 $d = 100\,\text{m}$를 걷는 데
$E = 49{,}000\,\text{J}$을 썼다고 하자. $g = 9.8\,\text{m/s}^2$을 쓰면

\begin{align*}
&m\,g\,d = 50 \times 9.8 \times 100 = 49{,}000 && \text{분모 (무게 $\times$ 거리)}\\
&\mathrm{CoT} = \dfrac{49{,}000}{49{,}000} = 1.0 && \text{무차원}
\end{align*}

CoT $=1.0$은 "몸무게 1\,N을 1\,m 옮기는 데 1\,J을 썼다"는 뜻이다. 같은 로봇이
$v = 1\,\text{m/s}$로 걸었다면 100\,m를 가는 데 100초가 걸리므로

\[
P = \frac{E}{t} = \frac{49{,}000}{100} = 490\,\text{W}, \qquad
\mathrm{CoT} = \frac{P}{m g v} = \frac{490}{50 \times 9.8 \times 1} = 1.0
\]

두 식이 같은 CoT를 주는 것은 \Cref{eq:lec02-cot}의 두 표현이 애초에 시간으로
약분한 관계이기 때문이다. 참고로 실제 이족보행 로봇의 CoT는 지형과 설계에 따라
1보다 훨씬 큰 경우가 많고, 바퀴 차량은 대체로 0.1 근처로 훨씬 작다. 이 비교값은
저자가 직관을 위해 임의로 지정한 것이며 특정 로봇의 실측값이 아니다.

\subsection{이동 개념의 특성화}
\label{sec:lec02-characterization}

이동 방식을 설계하거나 비교하려면 "무엇을 기준으로 좋고 나쁨을 말할 것인가"가
먼저 정해져야 한다. 슬라이드는 이 기준을 안정성, 접촉 특성, 환경 종류 세
갈래로 나눈다(\Cref{tab:lec02-characterization}).
% src: Lecture2-Locomotion.pdf slide 6

\begin{table}[ht]
\centering
\caption{이동 개념(locomotion concept)을 특성화하는 세 갈래 기준.}
\label{tab:lec02-characterization}
\begin{tabularx}{\linewidth}{@{}lL@{}}
\toprule
갈래 & 세부 항목 \\
\midrule
안정성(stability) & 접촉점 개수, 무게중심(COG) 위치, 정적/동적 안정화
  (static/dynamic stabilization), 지형의 경사(inclination of terrain) \\
접촉 특성(contact characteristics) & 점 접촉인지 면 접촉인지, 접촉각(angle of
  contact), 마찰(friction) \\
환경 종류(type of environment) & 구조화 정도(structure), 매질(medium: 물,
  공기, 무른 지면, 단단한 지면) \\
\bottomrule
\end{tabularx}
\end{table}

이 세 갈래는 \Cref{sec:lec02-walking-robots}의 다리 개수·관절 논의와
\Cref{sec:lec02-flying}의 비행 방식 비교에서 되풀이해 등장하는 잣대다.

\subsection{위치와 방향의 표현}
\label{sec:lec02-rotation}

\paragraph{In brief.}
\begin{description}
  \item[Why] 로봇의 자세(pose)를 계산하려면 위치뿐 아니라 방향(orientation)도
        숫자로 표현해야 한다. 관절을 하나씩 돌린 뒤 말단이 어디를 향하는지
        계산하는 모든 문제(정기구학 등)가 이 표현 위에서 돌아간다.
  \item[When] 좌표계 하나 안에서 벡터를 회전시키거나, 로봇 몸체 좌표계와
        전역(global) 좌표계 사이를 오갈 때.
  \item[Where] 이 절은 \Cref{ch:hw2}의 회전 행렬 배경(\Cref{sec:hw2-bg})과
        같은 수학을 쓴다. 다만 \Cref{ch:hw2}가 과제 문제 풀이에 집중했다면, 이
        절은 그 행렬이 이동(locomotion) 문맥에서 등장하는 자리(로봇의 자세
        계산)를 보인다.
  \item[How] 위치는 변위 벡터(displacement vector)로, 방향은 회전
        행렬(rotation matrix)로 표현하고, 여러 축을 돌릴 때는 행렬을 곱해
        합성한다.
\end{description}

\subsubsection{위치와 방향}
\label{sec:lec02-rotation-basic}

로봇의 자세(pose)는 두 가지로 나뉜다. \emph{위치(position)}는 한 좌표계의
원점에서 로봇까지의 변위 벡터(displacement vector) $(x, y, z)$로 표현하고,
\emph{방향(orientation)}은 좌표축이 서로 얼마나 돌아가 있는지를 나타내는 회전
행렬(rotation matrix) $R$로 표현한다.
% src: Lecture2-Locomotion.pdf slide 7

축 하나를 중심으로 벡터 $[x\ y\ z]$를 반시계 방향(counter-clockwise, CCW)으로
각도 $\theta$만큼 돌리는 기본 회전 행렬은 세 개다.
% src: Lecture2-Locomotion.pdf slide 18

\begin{equation}
\label{eq:lec02-rx}
R_X(\theta) =
\begin{bmatrix}
1 & 0 & 0 \\
0 & \cos\theta & -\sin\theta \\
0 & \sin\theta & \cos\theta
\end{bmatrix},
\qquad
R_Y(\theta) =
\begin{bmatrix}
\cos\theta & 0 & \sin\theta \\
0 & 1 & 0 \\
-\sin\theta & 0 & \cos\theta
\end{bmatrix}.
\end{equation}

\begin{equation}
\label{eq:lec02-rz}
R_Z(\theta) =
\begin{bmatrix}
\cos\theta & -\sin\theta & 0 \\
\sin\theta & \cos\theta & 0 \\
0 & 0 & 1
\end{bmatrix}.
\end{equation}

이 세 행렬은 \Cref{ch:hw2}의 \Cref{eq:hw2-rx,eq:hw2-ry,eq:hw2-rz}와 정확히
같은 행렬이다(같은 유도, \Cref{sec:hw2-bg-3d} 참고).

\paragraph{경고: 두 관례를 섞으면 부호가 뒤집힌다.}
\Cref{eq:lec02-rz}의 $R_Z(\theta)$는 오른쪽 위 성분이 $-\sin\theta$이고,
\emph{벡터를} 돌리는 능동적(active) 관례다(\Cref{sec:hw2-bg-two} 참고).
반면 \Cref{ch:lec01}의 $R(\theta)$는 오른쪽 위 성분이 $+\sin\theta$이고, 전역
좌표계를 로봇 좌표계로 옮기는 수동적(passive) 관례였다. 두 행렬은 서로
전치(transpose) 관계, 즉 $R_Z(\theta) = R(\theta)^{\mathsf{T}}$다. 어느 쪽을
쓰는지 밝히지 않고 두 장의 행렬을 섞어 쓰면 부호가 뒤집힌 답이 나온다. 이
장과 \Cref{ch:hw2}는 모두 능동적 관례($R_Z$, \Cref{eq:lec02-rz})를 쓴다.

\subsubsection{합성: 순서가 답을 바꾼다}
\label{sec:lec02-rotation-compose}

두 축을 연달아 돌리는 경우를 보자. 슬라이드는 $Z$축으로 $45^\circ$ 돌린 뒤,
"새 $X$축"으로 다시 $45^\circ$ 돌리는 예를 들고, 합성 행렬을
$R_{XZ} = R_X R_Z$로 적은 뒤 $V_{\text{new}} = R_{XZ} V$로 적용한다.
% src: Lecture2-Locomotion.pdf slide 26-27

\begin{equation}
\label{eq:lec02-rxz}
R_{XZ} = R_X(\theta)\,R_Z(\theta), \qquad V_{\text{new}} = R_{XZ}\,V.
\end{equation}

\Cref{ch:hw2}의 합성 규칙(\Cref{sec:hw2-bg-euler}, \Cref{eq:hw2-order})을 따르면,
행렬을 곱했을 때 벡터에 \emph{먼저} 닿는 쪽(오른쪽에 있는 행렬)이 먼저 적용된
회전이다. \Cref{eq:lec02-rxz}에서 $R_Z$가 오른쪽에 있으므로, 이 식은
\textbf{고정된 전역 $Z$축으로 먼저 돌리고, 그다음 고정된 전역 $X$축으로 돌리는
(외재적, extrinsic) 합성}이다.

\paragraph{경고: 슬라이드의 "새 $X$축"이라는 표현.}
슬라이드 캡션은 "새 $X$축(the new $X$ axis)"이라고 적었는데, 이는 보통 몸체와
함께 돌아간 축을 기준으로 다시 돈다는 내재적(intrinsic) 합성을 뜻하는 말이다.
그런데 실제로 적힌 식 $R_{XZ} = R_X R_Z$는 위에서 확인했듯 전역 축 기준의
외재적 합성이다. 만약 캡션 그대로 몸체의 새 $X$축을 기준으로 돌리려 했다면
합성 순서는 반대인 $R_Z R_X$가 되어야 한다(\Cref{sec:hw2-bg-euler} 참고, 외재적
"$A$ 다음 $B$"는 내재적 "$B$ 다음 $A$"와 같은 행렬이 된다). 이 불일치는
슬라이드 원문 자체에 있는 것으로, 이 장은 슬라이드에 \emph{적힌 식}
($R_{XZ} = R_X R_Z$, 전역 축 기준 "$Z$ 다음 $X$")을 그대로 채택하고 여기 기록만
남긴다.

핵심은 어느 쪽이 맞고 틀리냐가 아니라, \textbf{회전은 교환법칙이 성립하지
않으므로}($R_A R_B \neq R_B R_A$, \Cref{sec:hw2-bg-euler}) 순서를 반드시 함께
명시해야 한다는 점이다. 다음 예로 직접 확인한다.

\paragraph{예제: 순서를 바꾸면 다른 벡터가 나온다.}
$v_1 = (1, 2, 3)^{\mathsf{T}}$에 $\theta = 45^\circ$
($\cos 45^\circ = \sin 45^\circ = \sqrt{2}/2 \approx 0.7071$)로 $R_X$와 $R_Z$를
두 가지 순서로 합성한다.

\begin{align*}
R_Z(45^\circ)\,v_1 &= \left(-\tfrac{\sqrt{2}}{2},\ \tfrac{3\sqrt{2}}{2},\ 3\right)
   \approx (-0.71,\ 2.12,\ 3.00) && \text{$Z$ 먼저}\\
R_X(45^\circ)\bigl(R_Z(45^\circ)\,v_1\bigr)
   &\approx (-0.71,\ -0.62,\ 3.62) && \text{그다음 $X$ ($=R_{XZ}v_1$)}\\[6pt]
R_X(45^\circ)\,v_1 &\approx (1.00,\ -0.71,\ 3.54) && \text{$X$ 먼저}\\
R_Z(45^\circ)\bigl(R_X(45^\circ)\,v_1\bigr)
   &\approx (1.21,\ 0.21,\ 3.54) && \text{그다음 $Z$ ($=R_{ZX}v_1$)}
\end{align*}

두 결과 $(-0.71, -0.62, 3.62)$와 $(1.21, 0.21, 3.54)$는 서로 다른 벡터다.
회전은 길이를 보존하므로(\Cref{sec:hw2-bg-props}) 두 값 모두
$\lVert v \rVert^2 \approx 0.5 + 0.39 + 13.12 \approx 14.0$과
$1.46 + 0.04 + 12.53 \approx 14.0$으로 $\lVert v_1 \rVert^2 = 1+4+9=14$와
일치해 계산 자체는 맞지만, 벡터의 \emph{방향}은 순서에 따라 완전히 다르다.

\subsubsection{예제: 벡터를 직접 돌려 본다}
\label{sec:lec02-rotation-examples}

\paragraph{예제: $X$축으로 $90^\circ$.}
$v = (1, 2, 3)^{\mathsf{T}}$를 $X$축으로 $90^\circ$ CCW 돌린다
($\cos 90^\circ = 0$, $\sin 90^\circ = 1$).
% src: Lecture2-Locomotion.pdf slide 21

\begin{align*}
v' &= R_X(90^\circ)\,v
   = \begin{bmatrix} 1 & 0 & 0 \\ 0 & 0 & -1 \\ 0 & 1 & 0 \end{bmatrix}
     \begin{bmatrix} 1 \\ 2 \\ 3 \end{bmatrix}
   = \begin{bmatrix} 1 \\ -3 \\ 2 \end{bmatrix}
   && \text{슬라이드 답과 일치}
\end{align*}

\paragraph{예제: $Z$축으로 $45^\circ$.}
같은 $v = (1, 2, 3)^{\mathsf{T}}$를 $Z$축으로 $45^\circ$ CCW 돌린다.
% src: Lecture2-Locomotion.pdf slide 23

\begin{align*}
v' &= R_Z(45^\circ)\,v
   = \begin{bmatrix}
       \tfrac{\sqrt2}{2} & -\tfrac{\sqrt2}{2} & 0 \\
       \tfrac{\sqrt2}{2} & \tfrac{\sqrt2}{2} & 0 \\
       0 & 0 & 1
     \end{bmatrix}
     \begin{bmatrix} 1 \\ 2 \\ 3 \end{bmatrix}
   = \begin{bmatrix} -\tfrac{1}{\sqrt2} \\ \tfrac{3}{\sqrt2} \\ 3 \end{bmatrix}
   && \text{슬라이드 답과 일치}
\end{align*}

\paragraph{예제: 슬라이드의 두 번째 예제, 값을 다시 계산해 확인.}
점 $\langle 1, 1, 2\rangle$을 $\theta = 45^\circ$로 $Z$축 회전한 결과를 물었다.
% src: Lecture2-Locomotion.pdf slide 17

\begin{align*}
v' = R_Z(45^\circ)\begin{bmatrix}1\\1\\2\end{bmatrix}
   = \begin{bmatrix}
       \tfrac{\sqrt2}{2}(1) - \tfrac{\sqrt2}{2}(1) \\
       \tfrac{\sqrt2}{2}(1) + \tfrac{\sqrt2}{2}(1) \\
       2
     \end{bmatrix}
   = \begin{bmatrix} 0 \\ \sqrt2 \\ 2 \end{bmatrix}
   \approx \begin{bmatrix} 0 \\ 1.41 \\ 2 \end{bmatrix}
\end{align*}

pdftotext로 추출한 원문 슬라이드는 이 결과를 "0, 2, 2"로 보여 준다. 위 계산의
첫째·셋째 성분(0과 2)은 원문과 정확히 일치하고, 둘째 성분만 다르다. 근호
기호($\sqrt{\ }$)가 PDF 텍스트 추출 과정에서 자주 깨진다는 점(과제 설명에 명시된
알려진 문제)을 고려하면, 원문의 "2"는 $\sqrt{2}$에서 근호가 빠진 결과로 보는
것이 가장 설명이 잘 된다. 이 장은 직접 계산한 $(0,\ \sqrt2,\ 2)$를 최종 값으로
채택한다.

\subsection{다리 로봇}
\label{sec:lec02-walking-robots}

\paragraph{In brief.}
\begin{description}
  \item[Why] 바퀴가 지나가지 못하는 지형(구멍, 계단, 잔해)에서도 이동해야 할 때
        다리가 필요하다.
  \item[When] 지형이 불규칙하거나 사람이 쓰도록 만들어진 환경(계단, 좁은 통로)을
        지나야 할 때.
  \item[Where] 다리 하나당 자유도가 늘어날수록 적응력은 커지지만 에너지·질량·제어
        비용도 함께 커진다.
  \item[How] 다리마다 최소 2자유도(들어 올리기 + 휘두르기)를 주고, 걸음새(gait)로
        여러 다리의 동작을 맞춘다.
\end{description}

\paragraph{왜 다리가 필요한가.}
다리 달린 로봇은 거친 지형에서 적응력과 기동성이 뛰어나다. 바퀴로는 건널 수
없는 구멍이나 틈을 넘을 수 있고, 다리를 팔처럼 써서 물체를 조작할 잠재력도
있다. 하지만 그 대가로 소비 전력과 기계적 복잡도가 크게 늘어난다: 여러 자유도를
동시에 협조시켜 제어해야 하고, 발을 어디에 디딜지 정하려면 지형의 미세한 굴곡까지
센서로 봐야 한다.
% src: Lecture2-Locomotion.pdf slide 30

\paragraph{다리 하나에 관절이 몇 개 필요한가.}
다리를 앞으로 내밀려면 최소 2자유도가 필요하다: 들어 올리는 동작(lift)과 휘두르는
동작(swing)이다. 실제로는 3자유도를 쓰는 경우가 가장 흔하고, 발목에 4번째
자유도를 추가하면 걷기와 안정성이 좋아질 수 있다. 사람의 다리는 발가락 구동까지
합치면 7개가 넘는 주요 자유도를 쓴다.
% src: Lecture2-Locomotion.pdf slide 31-32
자유도를 하나 늘릴 때마다 설계와 제어의 복잡도가 늘어나고, 그 대가는 에너지
소비, 제어 난이도, 질량 증가로 돌아온다.

\subsection{걸음새(gait)의 수}
\label{sec:lec02-gaits}

\paragraph{고치는 문제.}
다리가 $k$개인 로봇은 각 다리를 들고 내리는 순서를 자유롭게 고를 수 있다. 이
순서의 가짓수, 즉 서로 다른 걸음새(distinct event sequence)의 개수를 세어보면
다리가 늘어날수록 그 수가 어떻게 되는지 감을 잡을 수 있다.

\paragraph{아이디어.}
다리 하나는 사건을 두 개 만든다: 드는 사건(lift)과 내려놓는 사건(release,
touch-down). 다리가 $k$개면 사건은 모두 $2k$개다. 그런데 걸음새는 반복되는
주기(cycle)이므로 사건 하나를 기준점(t=0)으로 고정할 수 있고, 남는 $2k-1$개
사건의 순서만 세면 된다. 그 순서의 가짓수가 곧 $(2k-1)!$이다.
% src: Lecture2-Locomotion.pdf slide 33 (숫자 2의 유래는 슬라이드 원문; "기준점 고정" 설명은 저자의 풀이)

\begin{equation}
\label{eq:lec02-gaits}
N = (2k - 1)!.
\end{equation}

\begin{align*}
&k=2 \ (\text{이족보행}): \quad N = (2\times2-1)! = 3! = 3\times2\times1 = 6 \\
&k=6 \ (\text{육족보행}): \quad N = (2\times6-1)! = 11! = 39{,}916{,}800
\end{align*}
% src: Lecture2-Locomotion.pdf slide 33

다리가 두 개에서 여섯 개로 늘었을 뿐인데 가짓수는 6에서 약 4천만으로 뛴다.
factorial(계승) 함수의 성질상 $k$가 조금만 커져도 $N$이 폭발적으로 커지기
때문이다. 이 폭발은 실제로 쓸 수 있는 걸음새를 사람이 하나하나 나열해서 고를 수
없다는 뜻이고, 그래서 실제 로봇은 이 조합 전체가 아니라 몇 가지 정형화된
걸음새(삼각보, 파상보 등)만 골라 쓴다.

\begin{table}[ht]
\centering
\caption{이족보행($k=2$)의 여섯 가지 사건 순서. R은 오른쪽 다리, L은 왼쪽
  다리를 뜻한다.}
\label{tab:lec02-biped-gaits}
\begin{tabularx}{\linewidth}{@{}cL@{}}
\toprule
번호 & 사건 순서 \\
\midrule
1 & 양다리 지지 $\to$ R 지지 / L 들림 $\to$ 양다리 지지 \\
2 & 양다리 지지 $\to$ R 들림 / L 지지 $\to$ 양다리 지지 \\
3 & R만 지지 $\to$ 양다리 들림 $\to$ 양다리 지지 \\
4 & R 지지 / L 들림 $\to$ R 들림 / L 지지 $\to$ R 지지 / L 들림 \\
5 & R 지지 / L 들림 $\to$ 양다리 들림 $\to$ R 지지 / L 들림 \\
6 & R 들림 / L 지지 $\to$ 양다리 들림 $\to$ R 들림 / L 지지 \\
\bottomrule
\end{tabularx}
\end{table}
% src: Lecture2-Locomotion.pdf slide 34

\Cref{tab:lec02-biped-gaits}는 \Cref{eq:lec02-gaits}가 준 6가지 조합을 슬라이드
그대로 나열한 것이다. 여섯 다리(육족) 로봇의 정적 걸음새(static gait)는 항상 세
다리로 이루어진 삼각대(tripod)가 지면에 닿아 있도록 짠다. 반대로 네 다리
동물에서 가장 자연스러운 걸음새(trot, gallop 등)는 대부분 동적(dynamic)
걸음새다.
% src: Lecture2-Locomotion.pdf slide 35-36

\subsection{정적 보행과 동적 보행}
\label{sec:lec02-static-dynamic}

\Cref{sec:lec02-bg-prereq}에서 정의한 정적/동적 안정성을 걸음새에 적용하면
\Cref{tab:lec02-static-dynamic}처럼 대비된다.
% src: Lecture2-Locomotion.pdf slide 38

\begin{table}[ht]
\centering
\caption{정적 보행(static walking)과 동적 보행(dynamic walking)의 대비.}
\label{tab:lec02-static-dynamic}
\begin{tabularx}{\linewidth}{@{}LL@{}}
\toprule
정적 보행 & 동적 보행 \\
\midrule
최소 세 다리가 항상 지면을 지지 & 지면 접촉 다리가 세 개 미만인 순간이 있음 \\
모든 관절이 그 순간 멈춰도 넘어지지 않음 & 계속 움직이지 않으면 넘어짐 \\
안전하지만 느리고 비효율적 & 빠르고 효율적이지만 구동·제어 부담이 큼 \\
\bottomrule
\end{tabularx}
\end{table}

\subsection{이족보행과 역진자 모형}
\label{sec:lec02-biped}

한 걸음은 사실 다각형을 굴리는 동작과 비슷하다. 다각형의 한 변의 길이가 보폭
(step length)에 해당하고, 보폭이 작아질수록 그 다각형은 원, 즉 바퀴에
가까워진다. 다만 자연은 회전 관절을 발명하지 않았고, 걷는 한 여전히 중력을
거스르는 일을 해야 한다.
% src: Lecture2-Locomotion.pdf slide 39
그래서 이족보행은 흔히 역진자(inverted pendulum, 받침점 위에서 거꾸로 서서
흔들리는 진자)로 모형화한다. 지지하는 다리가 받침점 역할을 하고, 몸통 COG가
진자의 추처럼 그 위에서 앞뒤로 넘어가듯 움직인다.
% src: Lecture2-Locomotion.pdf slide 40

\begin{figure}[ht]
\centering
\begin{tikzpicture}[scale=1.0,>=stealth]
  \draw[black!30] (-0.3,0) -- (5.3,0);
  \foreach \x in {0,1,2,3,4,5} {
    \draw[black!40] (\x,-0.05) -- (\x,0.05);
  }
  \draw[very thick] (0,0) -- (1,0.9);
  \draw[very thick] (1,0.9) -- (2,0);
  \draw[very thick,dashed,black!50] (1,0) -- (1,0.9);
  \filldraw (1,0.9) circle (2.2pt) node[above] {COG};
  \filldraw (0,0) circle (1.6pt);
  \filldraw (2,0) circle (1.6pt);
  \draw[<->,black!60] (0,-0.35) -- (2,-0.35)
        node[midway,below,minimum width=2.4cm] {보폭 (step length)};
  \draw[very thick] (3,0) -- (3.7,0.7);
  \draw[very thick] (3.7,0.7) -- (4.4,0);
  \draw[very thick,dashed,black!50] (3.7,0) -- (3.7,0.7);
  \filldraw (3.7,0.7) circle (2.2pt);
  \draw[very thick] (4.4,0) -- (4.9,0.35);
  \draw[very thick,dashed,black!50] (4.9,0) -- (4.9,0.35);
  \filldraw (4.9,0.35) circle (2.2pt);
  \filldraw (4.4,0) circle (1.6pt);
  \node[minimum width=2.6cm] at (3.9,-0.55) {보폭이 작을수록 원(바퀴)에 가까워짐};
\end{tikzpicture}
\caption{이족보행의 역진자 모형. 지지하는 다리가 받침점, 몸통 COG가 진자의
  추 역할을 한다. 보폭이 작을수록 걸음의 궤적은 다각형에서 원에 가까워진다.}
\label{fig:lec02-inverted-pendulum}
\end{figure}

탄성 에너지 저장(elastic energy storage, 힘줄이나 스프링처럼 착지 시 에너지를
저장했다가 다음 걸음에 되돌려주는 방식)을 활용하면 걸음새가 훨씬 효율적이
된다. 어떤 걸음새가 가장 경제적인지는 고정된 답이 있는 것이 아니라 원하는
속도(desired speed)에 따라 달라진다.
% src: Lecture2-Locomotion.pdf slide 41

\subsection{다리 로봇의 역사, 짧게}
\label{sec:lec02-history}

2000년대 이후 대표적인 이족보행 인간형 로봇은 HONDA ASIMO다. 2015년 DARPA
Robotics Challenge에서 여러 팀의 로봇이 문 열기, 계단 오르기 같은 과제에서
넘어지는 모습을 보였는데, 여기서 얻은 교훈은 "보행은 여전히 어렵다"는 것이었다.
이후 Boston Dynamics의 Spot(사족보행)과 Atlas(이족보행)가 훨씬 인상적인
기동성을 보여 주었다.
% src: Lecture2-Locomotion.pdf slide 42-46

\subsection{비행 로봇}
\label{sec:lec02-flying}

\paragraph{In brief.}
\begin{description}
  \item[Why] 지면이나 물이 아니라 공중에서 이동해야 할 때, 즉 지형 자체를
        무시하고 싶을 때 쓴다.
  \item[When] 넓은 영역을 빠르게 살피거나, 지면 접근이 아예 불가능한 상황.
  \item[Where] 비행 시간, 기동성, 기계적 복잡도가 방식마다 크게 다르므로 임무
        요구사항에 맞춰 골라야 한다.
  \item[How] 헬리콥터, 고정익, 비행선, 날갯짓 네 방식 중 하나를 고르고, 회전자
        속도 차이로 자세(roll, pitch, yaw)를 제어한다.
\end{description}

\Cref{tab:lec02-flight}은 무인 항공기(unmanned aerial vehicle, UAV)의 네 가지
비행 방식을 비교한다.
% src: Lecture2-Locomotion.pdf slide 46

\begin{table}[ht]
\centering
\caption{비행 개념(flight concept) 비교.}
\label{tab:lec02-flight}
\begin{tabularx}{\linewidth}{@{}lLL@{}}
\toprule
방식 & 비행 시간 및 특징 & 제약 \\
\midrule
헬리콥터 & 20분 미만 & 매우 동적이고 기동성이 높음 \\
고정익(fixed wing) 항공기 & 수 시간, 연속 비행 가능 & 비홀로노믹
  (non-holonomic) 제약 \\
비행선(blimp, 공기보다 가벼움) & 수 시간(바람 조건에 의존) & 바람에 민감,
  탑재 중량 대비 크기가 큼 \\
날갯짓(flapping wings) & 20분 미만, 활공(gliding) 가능 & 비홀로노믹 제약,
  기구가 매우 복잡함 \\
\bottomrule
\end{tabularx}
\end{table}

여기서 비홀로노믹(non-holonomic)이란 "그 순간 갈 수 있는 방향이 원하는 방향
전부가 아니라 일부로 제한된다"는 뜻이다. 예를 들어 고정익 비행기는 옆으로 미끄러져
이동할 수 없고 앞으로 나아가면서 방향을 서서히 바꿔야 한다. 자동차가 옆으로 바로
평행이동하지 못하는 것과 같은 종류의 제약이다.

\paragraph{쿼드콥터(quadcopter)의 자세 제어.}
쿼드콥터는 회전자(rotor) 네 개의 회전 속도 차이만으로 세 가지 회전, 즉
좌우로 기울이는 롤(roll), 앞뒤로 기울이는 피치(pitch), 수직축으로 도는
요(yaw)를 만든다.
% src: Lecture2-Locomotion.pdf slide 47

\begin{figure}[ht]
\centering
\begin{tikzpicture}[scale=1.0,>=stealth]
  % body cross
  \draw[very thick] (-1.6,0) -- (1.6,0);
  \draw[very thick] (0,-1.6) -- (0,1.6);
  \foreach \p in {(-1.6,0),(1.6,0),(0,-1.6),(0,1.6)} {
    \filldraw[black!70] \p circle (5pt);
  }
  \draw[->,red!70!black,very thick] (0,0) -- (2.1,0)
        node[right,minimum width=1.6cm] {Roll ($x$)};
  \draw[->,blue!70!black,very thick] (0,0) -- (0,2.1)
        node[above,minimum width=1.6cm] {Pitch ($y$)};
  \draw[->,black!70,very thick] (0,0) circle (0.55);
  \node[minimum width=1.6cm] at (0.95,-0.95) {Yaw ($z$, 위에서 본 회전)};
\end{tikzpicture}
\caption{쿼드콥터의 세 회전축. Roll과 Pitch는 마주보는 회전자 한 쌍의 속도
  차이로, Yaw는 시계 방향/반시계 방향으로 도는 두 쌍의 속도 차이로 만든다.}
\label{fig:lec02-quadcopter}
\end{figure}

마주 보는 회전자 두 개(예: 앞뒤 쌍)의 속도를 다르게 하면 그 축으로 몸체가
기울며 피치가 생기고, 다른 방향의 쌍을 다르게 하면 롤이 생긴다. 네 회전자는
시계 방향 두 개, 반시계 방향 두 개로 짝지어 회전 반작용 토크를 상쇄하는데, 이
두 쌍의 속도를 서로 다르게 하면 반작용 토크가 상쇄되지 않고 남아 기체가
수직축을 중심으로 요(yaw) 회전을 한다.

\subsection{수영 로봇}
\label{sec:lec02-swimming}

수중 이동 로봇의 예로는 로봇 물고기 무리(swarm of robotic fish)가 있다. 여러
마리가 무리를 지어 헤엄치며 환경을 탐사하는 방식으로, 슬라이드는 이 한 가지
사례만을 제시한다.
% src: Lecture2-Locomotion.pdf slide 56

\section{어떻게 적용하는가}
\label{sec:lec02-apply}

이 장의 회전 행렬 대수가 바로 \Cref{ch:hw2}의 내용이다. 그 과제는
$v_1 = (1, 2, 3)^{\mathsf{T}}$에 $R_X$, $R_Y$, $R_Z$를 $XYZ$ 순서
(\Cref{sec:hw2-q2})와 $YZX$ 순서(\Cref{sec:hw2-q3})로 합성해 서로 다른 결과
$v_2$, $v_3$을 얻는다(\Cref{tab:hw2-results}). 이때 쓰는 $R_X$, $R_Y$, $R_Z$는
바로 이 장의 \Cref{eq:lec02-rx,eq:lec02-rz}와 같은 능동적(active) 관례다.

아래 코드는 이 장의 $R_Z$를 만들어 전역 좌표계와 로봇 좌표계 사이에서 점 하나를
옮긴다. 전치를 취하면 \Cref{ch:lec01}의 관례로 바뀐다는 점을 주석으로 남겼다.

\begin{lstlisting}[style=py,caption={$R_Z$로 점을 전역/로봇 좌표계 사이에서 옮기기.}]
import numpy as np

def rz(theta):
    c, s = np.cos(theta), np.sin(theta)
    return np.array([[c, -s, 0],
                      [s,  c, 0],
                      [0,  0, 1]])

theta = np.pi / 4
R = rz(theta)                  # 이 장의 능동적 관례, +theta = CCW
p_robot = np.array([1.0, 1.0, 2.0])

p_global = R @ p_robot         # 로봇 좌표계 -> 전역 좌표계 (벡터를 돌린다)
p_back = R.T @ p_global        # R.T = R^-1 : 전역 -> 로봇 좌표계로 되돌리기
                                # R.T를 쓰는 순간 lec01의 수동적 관례와 같아진다
\end{lstlisting}

\Cref{ch:hw1}에서 다룬 시뮬레이터 Webots와 ARGoS는 모두 차동구동
(differential-drive) 로봇을 전역 좌표계와 로봇 좌표계 두 틀 모두로 조작할 수
있게 해 준다. 시뮬레이터의 API가 로봇 좌표계 기준 속도(전진 속도, 각속도)를
받으면, 화면에 그려지는 전역 좌표계의 위치·방향은 결국 이 절의 $R_Z$(또는
\Cref{ch:lec01}의 $R^{\mathsf{T}}_Z$)로 매 스텝 갱신된 결과다.

\section{앞으로}
\label{sec:lec02-future}

다리 이동(legged locomotion)은 바퀴 이동에 비해 아직 연구 단계에 머무는
부분이 많다. 걸음새 최적화(gait optimization)와 탄성 에너지 저장을 결합해
효율을 끌어올리는 문제, 속도와 에너지 효율 사이의 절충, 걷는 동안 실시간으로
균형을 잡는 동적 균형 제어(dynamic balance control), 발을 디딜 위치를 정하기
위한 정밀한 지형 인식(terrain perception)이 모두 열린 문제로 남아 있다. 반대로
바퀴 이동에서는 방향과 위치를 다루는 수학이 이미 정리되어 있으므로, 다음으로
이어지는 주제는 바퀴가 지켜야 하는 운동학적 제약, 즉 구름(rolling)과 옆으로
미끄러지지 않는다는(no lateral sliding) 제약이다. 이는 \Cref{ch:lec01}의
운동 제어(motion control) 절에서 이어받는다.

% Back to English prose - restore cleveref's default names (hw2.tex 206-216).
\crefname{equation}{Equation}{Equations}
\Crefname{equation}{Equation}{Equations}
\crefname{figure}{Figure}{Figures}
\Crefname{figure}{Figure}{Figures}
\crefname{table}{Table}{Tables}
\Crefname{table}{Table}{Tables}
\crefname{section}{Section}{Sections}
\Crefname{section}{Section}{Sections}
\crefname{chapter}{Chapter}{Chapters}
\Crefname{chapter}{Chapter}{Chapters}
